\section{Background and Related Work}\label{sec:related}

\textbf{Carbon-aware shifting.}
Temporal shifting defers flexible jobs to hours with lower grid carbon intensity (CI);
spatial shifting moves them to cleaner regions~\cite{wiesner2021wait,radovanovic2023carbon}.
Systems apply the idea to elastic batch work~\cite{hanafy2023carbonscaler,souza2023ecovisor},
cloud cost--carbon trade-offs~\cite{hanafy2024gaia}, data-centre design~\cite{acun2023carbonexplorer},
precedence-constrained data processing~\cite{lechowicz2025pcaps}, and GPU clusters for machine
learning, where GREEN~\cite{xu2025green} shifts training jobs toward low-carbon periods.
Online-algorithm work gives competitive guarantees for pause/resume and conversion
problems~\cite{lechowicz2023pauseresume,lechowicz2024conversion}.
Sukprasert et al.~\cite{sukprasert2024limits} bound the achievable benefit of shifting with
perfect information and find it smaller than often assumed; reviews survey spatio-temporal
schemes~\cite{asadov2025review,riepin2025truly}. These works mostly assume point forecasts or
perfectly known CI, and evaluate aggregate savings. None measures how many \emph{individual} jobs are made
worse off, nor separates forecast error from contention among deferred jobs.

\textbf{CI forecasting.}
CarbonCast and successors forecast CI days ahead from generation-mix and weather
features~\cite{maji2022carboncast,maji2023multiday,maji2022dacf,yan2025ensembleci}; CarbonX adds
prediction intervals~\cite{maji2025carbonx}. The GB system operator (NESO) publishes a 48\,h national
forecast~\cite{ciapi,bruce2024cimethod}, but its API keeps only the latest forecast per
half-hour~\cite{ciapi_issue11}; horizon-resolved history survives only in a third-party
archive~\cite{powell_citracking}, which we use.

\textbf{Uncertainty-aware scheduling.}
Conformal prediction gives distribution-free, finite-sample marginal coverage for
forecasts~\cite{angelopoulos2023gentle}; adaptive conformal inference (ACI) keeps long-run coverage
under distribution shift~\cite{gibbs2021aci}, and time-series variants exist~\cite{xu2023spci,xu2023cpts}.
Closest to us, Li et al.~\cite{li2024uncertainty} wrap CI forecasts in conformal intervals and show
that ignoring uncertainty can increase emissions in aggregate. That line of work treats
carbon-forecast uncertainty as the primary obstacle to deferral. Our decomposition shows that forecast
error accounts for \res{V2U}{POOL-S12-utiltarget0.5}{P2}{shareforecast}\% of harmful excess emissions
at 50\% utilisation and only \res{V2U}{POOL-S12-utiltarget0.9}{P2}{shareforecast}\% at 90\%. Contention
among deferred jobs dominates at realistic loads. Others use robust or
chance-constrained optimisation~\cite{hall2025probabilistic,yang2025conformalccp,ruparel2026dro} or
uncertainty-quantified advice in online algorithms~\cite{johnson2025uqadvice}.
We differ in asking \emph{why} deferral backfires at the level of individual jobs on a GPU trace,
in attributing the backfire to forecast error, re-planning and contention, and in treating
uncertainty-aware rules as points on a savings--harm trade-off rather than as a fix.
